\section*{Chapitre \ref{ch:g8:reproductive-systems} --- Les appareils reproducteurs}

\begin{solution}{exo:g8:reproductive-systems:1}
Les \omterm{def:g8:reproductive-systems:male}{testicules}, dans les bourses~: production continue de \omterm{def:g8:reproductive-systems:male}{cellules
mâles} à partir de la \omterm{def:g8:puberty:puberty}{puberté}. Les conduits et les glandes~: le
stockage, le transport, et les liquides qui font le sperme. Le
\omterm{def:g8:reproductive-systems:male}{pénis}~: la sortie du sperme --- et celle de l'\omterm{def:g7:removing-wastes:kidneys}{urine}, dans sa
fonction séparée et gardée par des \omterm{def:g7:heart-and-circulation:rooms}{valves}.
\end{solution}

\begin{solution}{exo:g8:reproductive-systems:2}
Les \omterm{def:g8:reproductive-systems:female}{ovaires}~: le stock de \omterm{def:g4:animal-reproduction:sexual}{cellules œufs} depuis la \omterm{def:g2:born-grow-die:cycle}{naissance}, dont
une est libérée par cycle à partir de la \omterm{def:g8:puberty:puberty}{puberté}. Les \omterm{def:g8:reproductive-systems:female}{trompes}~:
elles recueillent la \omterm{def:g4:animal-reproduction:sexual}{cellule œuf} libérée --- le lieu habituel de la
\omterm{def:g5:flower-to-fruit:steps}{fécondation}. L'\omterm{prop:g5:human-reproduction:plan}{utérus}~: le logement musclé, à muqueuse renouvelée.
Le vagin~: la réception du sperme, et la voie de la \omterm{def:g2:born-grow-die:cycle}{naissance}.
\end{solution}

\begin{solution}{exo:g8:reproductive-systems:3}
\omterm{def:g8:reproductive-systems:male}{Cellule mâle}~: la plus petite du corps --- une tête pleine de \omterm{def:g6:cell-unit-of-life:parts}{noyau}
et une queue de nage, peu pourvue, fabriquée en continu par
centaines de millions. \omterm{def:g4:animal-reproduction:sexual}{Cellule œuf}~: la plus grande du corps, ronde
et immobile, richement pourvue, libérée une par mois.
\end{solution}

\begin{solution}{exo:g8:reproductive-systems:4}
Un \omterm{def:g8:reproductive-systems:female}{ovaire} libère une \omterm{def:g4:animal-reproduction:sexual}{cellule œuf}, recueillie par sa \omterm{def:g8:reproductive-systems:female}{trompe} --- vers
le jour 14. Elle peut être fécondée pendant environ une journée.
\end{solution}

\begin{solution}{exo:g8:reproductive-systems:5}
La destruction et l'évacuation de la muqueuse préparée de l'\omterm{prop:g5:human-reproduction:plan}{utérus},
qui s'écoule pendant quelques jours. Leur arrivée dit que ce cycle
s'est terminé sans \omterm{def:g5:flower-to-fruit:steps}{fécondation}.
\end{solution}

\begin{solution}{exo:g8:reproductive-systems:6}
La \omterm{def:g5:flower-to-fruit:steps}{fécondation}. Le \omterm{prop:g8:sexual-reproduction:core}{zygote} s'installe dans la muqueuse préparée, le
cycle s'interrompt, et la \omterm{def:g5:human-reproduction:pregnancy}{grossesse} commence.
\end{solution}

\begin{solution}{exo:g8:reproductive-systems:7}
Étape 1, les \omterm{def:g8:sexual-reproduction:gametes}{gamètes}~: \omterm{def:g8:reproductive-systems:male}{testicules} et \omterm{def:g8:reproductive-systems:female}{ovaires}. Étape 2, la
\omterm{def:g5:flower-to-fruit:steps}{fécondation}~: dans la \omterm{def:g8:reproductive-systems:female}{trompe}, atteinte par la nage des \omterm{def:g8:reproductive-systems:male}{cellules
mâles} depuis le vagin à travers l'\omterm{prop:g5:human-reproduction:plan}{utérus}. Étape 3, le développement
du \omterm{prop:g8:sexual-reproduction:core}{zygote} qui se divise~: logé dans la muqueuse préparée de
l'\omterm{prop:g5:human-reproduction:plan}{utérus}.
\end{solution}

\begin{solution}{exo:g8:reproductive-systems:8}
Le côté masculin suit la logique du «~beaucoup et bon marché~»~: des
millions de nageuses en continu, chacune sacrifiable. Le côté
féminin suit celle du «~peu et bien pourvu~»~: une seule \omterm{def:g4:animal-reproduction:sexual}{cellule œuf}
garnie par mois, et derrière elle tout l'investissement du logement
utérin. La ligne des stratégies, écrite en \omterm{def:g5:first-look-at-cells:cell}{cellules}.
\end{solution}

\begin{solution}{exo:g8:reproductive-systems:9}
La machinerie se règle pendant ses premières années --- la chaîne
d'\omterm{def:g8:puberty:hormones}{hormones} trouve son rythme --- si bien que l'irrégularité est
attendue~; et la durée du cycle varie réellement, d'une personne à
l'autre et d'un mois à l'autre. Les deux sont l'appareil qui
fonctionne normalement.
\end{solution}

\begin{solution}{exo:g8:reproductive-systems:10}
L'\omterm{prop:g5:human-reproduction:plan}{utérus} est un \omterm{def:g3:bones-and-muscles:muscle}{muscle}, et l'évacuation est un travail musculaire
--- ses contractions se ressentent comme des douleurs. La chaleur et
le mouvement aident (et l'infirmière scolaire a des conseils quand
ils ne suffisent pas).
\end{solution}

\begin{solution}{exo:g8:reproductive-systems:11}
Le métier de la \omterm{def:g8:reproductive-systems:male}{cellule mâle} est le voyage~: bon marché, nageuse, et
nécessaire par millions parce que presque aucune n'arrive --- d'où
une usine qui tourne en continu. Le métier de la \omterm{def:g4:animal-reproduction:sexual}{cellule œuf} est de
devenir les premières provisions du \omterm{prop:g8:sexual-reproduction:core}{zygote}~: coûteuse, en stock
depuis la \omterm{def:g2:born-grow-die:cycle}{naissance}, et libérée avec un mois de préparation derrière
elle --- d'où une usine qui est un trésor, et qui dépense une unité
à la fois.
\end{solution}

\begin{solution}{exo:g8:reproductive-systems:12}
L'événement 1 est la question en préparation~: la muqueuse
reconstruite, le logement prêt. L'événement 2 la pose~: une \omterm{def:g4:animal-reproduction:sexual}{cellule
œuf} libérée --- la \omterm{def:g5:flower-to-fruit:steps}{fécondation} viendra-t-elle~? L'événement 3 est la
réponse «~non~»~: la muqueuse inutile s'évacue, et la roue repose la
question le mois suivant. L'événement 4 est la réponse «~oui~»~: le
\omterm{prop:g8:sexual-reproduction:core}{zygote} s'installe, la question cesse d'être posée, et le logement
sert son but pendant neuf mois.
\end{solution}

\begin{solution}{pb:g8:reproductive-systems:1}
\textbf{1.} Des \omterm{def:g8:reproductive-systems:male}{testicules} --- les usines, tenues dans les bourses
--- les \omterm{def:g8:reproductive-systems:male}{cellules mâles} passent dans les conduits pour y être
stockées et transportées, sont rejointes par les liquides des
glandes pour former le sperme, et sortent par le \omterm{def:g8:reproductive-systems:male}{pénis}~; le même
organe évacue l'\omterm{def:g7:removing-wastes:kidneys}{urine} à d'autres moments, un jeu de \omterm{def:g7:heart-and-circulation:rooms}{valves} tenant
les deux fonctions strictement séparées.

\textbf{2.} Le stock attend dans les \omterm{def:g8:reproductive-systems:female}{ovaires}, depuis la \omterm{def:g2:born-grow-die:cycle}{naissance}~;
la \omterm{def:g5:flower-to-fruit:steps}{fécondation} se produit habituellement dans les \omterm{def:g8:reproductive-systems:female}{trompes}~; une
\omterm{def:g5:human-reproduction:pregnancy}{grossesse} se loge dans la muqueuse préparée de l'\omterm{prop:g5:human-reproduction:plan}{utérus}.

\textbf{3.} Les \omterm{def:g8:sexual-reproduction:gametes}{gamètes} sont bâtis selon des stratégies opposées~:
des millions de nageuses bon marché contre une seule \omterm{def:g4:animal-reproduction:sexual}{cellule œuf}
bien pourvue par mois --- le «~nombre contre soin~», inscrit dans
les \omterm{def:g5:first-look-at-cells:cell}{cellules} elles-mêmes.

\textbf{4.} Ils attendaient leurs ordres~: les appareils étaient
bâtis mais à l'arrêt jusqu'à ce que les \omterm{def:g8:puberty:hormones}{hormones} du centre de
commande de la base du \omterm{def:g5:brain-in-command:system}{cerveau} les mettent en marche à la \omterm{def:g8:puberty:puberty}{puberté}.

\textbf{5.} Jours 1 à 5 environ~: les règles. Jusqu'au milieu du
cycle~: la muqueuse se reconstruit. Vers le jour 14~: l'\omterm{prop:g8:reproductive-systems:cycle}{ovulation},
avec le jour fécond autour d'elle. Puis l'attente préparée ---
jusqu'au jour 28 et au tour suivant de la roue.

\textbf{6.} La durée du cycle varie d'une personne à l'autre, et
d'un mois à l'autre chez une même personne --- et les premières
années sont irrégulières pendant que la machinerie se règle. Le
calendrier enseigne le schéma, pas les dates exactes de quiconque.

\textbf{7.} Annulés~: la destruction de la muqueuse et les règles
suivantes --- la roue s'arrête. Commencés~: l'installation du \omterm{prop:g8:sexual-reproduction:core}{zygote}
dans la muqueuse, et les neuf mois de la \omterm{def:g5:human-reproduction:pregnancy}{grossesse}.

\textbf{8.} Les douleurs sont l'\omterm{prop:g5:human-reproduction:plan}{utérus} musclé qui travaille pendant
l'évacuation --- normal, ce n'est pas une lésion~; la chaleur et le
mouvement aident, et tout ce qui est plus fort a son adresse à la
porte de l'infirmerie ou du cabinet médical.

\textbf{9.} Par exemple~: «~Ce sont des organes comme le \omterm{def:g4:heart-and-blood:heart}{cœur} et les
\omterm{def:g7:how-animals-breathe:four}{poumons} --- chacun ici possède l'un des deux jeux, et voici tout
simplement leur chapitre.~»

\textbf{10.} Usines~: \omterm{def:g8:reproductive-systems:male}{testicules} --- \omterm{def:g8:reproductive-systems:female}{ovaires}. \omterm{def:g8:sexual-reproduction:gametes}{Gamètes}~: des millions
de nageuses en continu --- une \omterm{def:g4:animal-reproduction:sexual}{cellule œuf} bien pourvue par mois.
Rencontre~: la \omterm{def:g8:reproductive-systems:female}{trompe}. Logement~: aucun --- l'\omterm{prop:g5:human-reproduction:plan}{utérus} préparé chaque
mois. Calendrier~: continu depuis la \omterm{def:g8:puberty:puberty}{puberté} --- cyclique, de la
\omterm{def:g8:puberty:puberty}{puberté} à cinquante ans environ.

\textbf{11.} Étape 1 sur les deux modèles~: les usines fabriquent
les \omterm{def:g8:sexual-reproduction:gametes}{gamètes}. Étape 2 sur le modèle féminin~: leur rencontre dans la
\omterm{def:g8:reproductive-systems:female}{trompe}. Étape 3, entre le modèle et le calendrier~: le \omterm{prop:g8:sexual-reproduction:core}{zygote} qui se
loge dans la muqueuse préparée du calendrier, la roue arrêtée, le
développement commencé.

\textbf{12.} «~Parce que les corps fonctionnent mieux avec du savoir
qu'avec des légendes --- et que tout ce qui a été enseigné
aujourd'hui a une adresse formée et confidentielle pour les
questions qui suivront.~»
\end{solution}
